\documentclass[reqno, 12pt]{amsart}

\usepackage{amssymb}
\usepackage{amsthm}
\usepackage{amsmath}
\usepackage{amsxtra}
\usepackage{mathrsfs}
\usepackage{comment}
\usepackage{graphicx}
\usepackage{placeins}
\usepackage{multicol}
\usepackage{bbm}
\usepackage{stmaryrd}
%\usepackage{enumerate}
\usepackage[inline, shortlabels]{enumitem}
\usepackage{mathdots}
\usepackage{colonequals}
\usepackage{hyperref}


\usepackage{calc}
\usepackage{tikz}
\usetikzlibrary{arrows,calc,automata,shadows,backgrounds,positioning,intersections,fadings,decorations.pathreplacing,shapes,matrix}
\usepackage{tikz-cd}

\usepackage{fullpage}

\newtheorem{thm}{Theorem}[section]
\newtheorem{lem}[thm]{Lemma}
\newtheorem{cor}[thm]{Corollary}
\newtheorem{conj}[thm]{Conjecture}
\newtheorem{prop}[thm]{Proposition}
\theoremstyle{definition}  
\newtheorem{defn} [thm] {Definition} 
\newtheorem{claim}[thm]{Claim}
\newtheorem{example} [thm] {Example}
\newtheorem{rem} [thm] {Remark}

\newenvironment{problab}[1]
{\noindent\textbf{Problem #1}.}
{\vskip 6pt}

\newenvironment{exolab}[1]
{\noindent\textbf{Exercise #1}.}
{\vskip 6pt}

\theoremstyle{remark}
\newtheorem*{soln}{Solution}

\newcommand{\C}{\mathbb C}
\renewcommand{\H}{\mathbb H}
\newcommand{\D}{\mathbb D}
\newcommand{\F}{\mathbb F}
\newcommand{\Q}{\mathbb Q}
\newcommand{\R}{\mathbb R}
\newcommand{\Z}{\mathbb Z}
\newcommand{\T}{\mathbb T}
\renewcommand{\P}{\mathcal P}
\renewcommand{\O}{\mathcal O}
\newcommand{\m}{\mathfrak m}
\newcommand{\A}{\mathbb A}
\renewcommand{\SS}{\mathbb S}
\renewcommand{\L}{\mathcal L}

\newcommand{\B}{\mathcal B}
\newcommand{\E}{\mathcal E}

\newcommand{\SSpan}{\text{span}}
\newcommand{\la}{\langle}
\newcommand{\ra}{\rangle}

\DeclareMathOperator{\lcm}{lcm}
\DeclareMathOperator{\img}{img}
\DeclareMathOperator{\Ann}{Ann}
\DeclareMathOperator{\Tor}{Tor}
\DeclareMathOperator{\tr}{tr}
\DeclareMathOperator{\Hom}{Hom}
\DeclareMathOperator{\End}{End}
\DeclareMathOperator{\Aut}{Aut}
\DeclareMathOperator{\Gal}{Gal}
\DeclareMathOperator{\sgn}{sgn}
\DeclareMathOperator{\ord}{ord}
\DeclareMathOperator{\ad}{ad}
\DeclareMathOperator{\coker}{coker}
\DeclareMathOperator{\rank}{rank}
\DeclareMathOperator{\minpoly}{minpoly}
\DeclareMathOperator{\charpoly}{charpoly}

\DeclareMathOperator{\id}{id}

\usepackage{mathpazo}

\everymath{\displaystyle}

\tikzset{commutative diagrams/.cd, arrow style = tikz, diagrams = {>=latex}}

\newenvironment{amatrix}[1]{%
  \left(\begin{array}{@{}*{#1}{c}|c@{}}
}{%
  \end{array}\right)
}

\usepackage{abstract}
\renewcommand{\abstractname}{}    % clear the title
\renewcommand{\absnamepos}{empty} % originally center
%\setlength{\abstitleskip}

\begin{document}
%\renewcommand{\abstractname}{\vspace{-\baselineskip}}

\title{18.700 Problem Set 9 (Optional)}
%\author{Évariste Galois}
\dedicatory{This is an optional problem set that is not to be turned in, and will not be graded.}
%\thanks{}
%\contrib{}
%\address{}

\maketitle

\begin{abstract}
  \noindent Collaborated with:\\
  Sources used: 
\end{abstract}
\vskip 0.5cm

%Let $V$ and $W$ be nonzero finite-dimensional inner product spaces over $\F$.
%\vskip 0.5cm

\begin{problab}{1} (8A \#2) %(7 points)
  Suppose $T \in \L(V)$, $m \in \Z_{>0}$, and $v \in V$ such that $T^{m-1}(v) \neq 0$ but $T^m(v) = 0$. Prove that $v, T(v), T^2(v), \ldots, T^{m-1}(v)$ is linearly independent.
\end{problab}

%\begin{soln}
%\end{soln}

\begin{problab}{2} (8A \#9) %(6 points)
  Suppose $T \in \L(V)$ and $m \in \Z_{\geq 0}$. Prove that \begin{equation*}
    \ker(T^m) = \ker(T^{m+1}) \iff \img(T^m) = \img(T^{m+1}) \, .
  \end{equation*}
\end{problab}

\begin{problab}{3} (8B \#9) %(7 points)
  Suppose $\F = \C$ and $T \in \L(V)$. Prove that there exist $D, N \in \L(V)$ such that $T = D + N$, the operator $D$ is diagonalizable, $N$ is nilpotent, and $DN = ND$.
\end{problab}

\begin{problab}{4} (8B \#16) %(6 points)
  Let $T$ be the operator on $\C^6$ defined by
  \begin{equation*}
    T(z_1, z_2, z_3, z_4, z_5, z_6) = (0, z_1, z_2, 0, z_4, 0) \, .
  \end{equation*}
  Find $\minpoly(T)$ and $\charpoly(T)$.
\end{problab}

\begin{problab}{5} (8C \#5) %(6 points)
  Suppose $T \in \L(\C^2)$ is the operator defined by
  \begin{equation*}
    T(w,z) = (-w - z, 9w + 5z) \, .
  \end{equation*}
  Find a Jordan basis for $T$.
\end{problab}

\begin{problab}{6} (8D \#5) %(4 points)
  Suppose $V$ is an inner product space. Suppose $T \in \L(V)$ is a positive operator and $\tr(T) = 0$. Prove that $T = 0$.
\end{problab}

\end{document} 	
